\section{Discussion}



\subsection{Transition to Agentic Market Microstructure}

Our analysis of the interaction dynamics between the platform and non-human actors suggests that the current static pricing models are insufficient for an agent-mediated economy. If we assume a transition toward a direct revelation mechanism, where actors must reveal their true valuation of a good through bidding dynamics, we inevitably introduce significant stochasticity into the pricing system. Unlike traditional e-commerce where prices are relatively sticky, such a mechanism implies a high volatility characteristic of financial equity markets (without the fungability however).

However, ecommerce commodities differ fundamentally from financial securities: they possess a hard floor defined by unit economics and reservation prices. The market might react enthusiastically to an iPhone priced at \$1, such a transaction is not permissible. The platform must establish an initial valuation anchor ($P_{0}$) defined by the marginal cost plus a target margin, around which the market price is permitted to fluctuate. We float the introduction of GenAI Agents as Institutional Market Makers. As the arms race for greater autonomy of agnetic systems grows, the commercial viability of AI agents has the potential to disseminate into every-day users directly interacting with them rather than e-commerce platforms. This is also under the assumption of expected transactional capabilities being given to AI Agents.



\subsection{Risk Assessment and Limitations}
\label{sec:limitations_risks}

This technology does not come without a more bitter side, ethical concerns do arise from the idea of deploying black-box like solutions to set prices based on a behavioral attributes. Approaches like universal behavioral profile modeling (UBPM) used in recommendation systems is very broadly utilized.

In our experimental setup we randomly assign each user to a platform and, within that platform, assign them to a task. Figure~\ref{fig:exp_design_tree} summarizes this design decision tree.

\begin{figure}[ht]
  \centering
  \resizebox{0.92\columnwidth}{!}{%
    % Horizontal tree: level distance must exceed ~half parent + half child width or nodes overlap (resizebox does not fix that).
\begin{tikzpicture}[
  grow=right,
  level distance=30mm,
  sibling distance=23mm,
  decision/.style={
    rectangle,
    draw,
    rounded corners=1.5pt,
    align=center,
    inner sep=1.2pt,
    minimum width=14mm,
    minimum height=4.8mm,
    font=\scriptsize,
  },
  leaf/.style={
    rectangle,
    draw,
    align=center,
    inner sep=1.2pt,
    text width=19mm,
    minimum height=4mm,
    font=\scriptsize,
  },
  edge from parent/.style={draw, -{Latex[length=1.2mm]}},
]
\node[decision] {Participant}
  child {
    node[decision] {Platform: Hotel}
      child {node[leaf] {Task sampled\\from hotel pool}}
  }
  child {
    node[decision] {Platform: Airline}
      child {node[leaf] {Task sampled\\from airline pool}}
  };
\end{tikzpicture}

  }
  \caption{Experimental design decision tree for participant assignment.}
  \label{fig:exp_design_tree}
\end{figure}

Although our participant sample size is somewhat low for humans, we do a one-to-one balance of human-to-agent experimental sessions. This way we are observing a uniform distribution of participation from each participating side. Our sample size of participants might look scarce, but each participant generates a rich amount of data, with a totality of 3,874 rows of data.

With a system like this there is potential for strong drift given the rapid advance of agentic systems and user preference. Our intent behind adding the UX term into the reward shaping process was to further address the risk of degraded user experience. Looking deeper at the underlying methodology, reinforcement learning does not come without it's complications such as reward hacking and often the lack of intepretability which is quite critical in systems that have a strong impact on the revenue of a company.

% \subsection{Implications of Findings} Interpretation of results and altenrative scenarios with broader market implications.
